\pagestyle{empty}
\tight
\begin{tblr}{width=\textwidth,colspec={|Q[c,m,1.9cm]|X[l,m]|},hlines,vlines,hspan=minimal,row{1}={bg=headblue},rowsep=1pt,colsep=3pt}
\SetCell{c,m}\hfil\logo\hfil & \textbf{POLITEKNIK NEGERI MALANG}\par\textbf{JURUSAN TEKNOLOGI INFORMASI}\par\textbf{PROGRAM STUDI: D4 TEKNIK INFORMATIKA} \\
\SetCell[c=2]{c} \textbf{RENCANA PEMBELAJARAN SEMESTER (RPS)} & \\
\end{tblr}
\begin{longtblr}{width=\textwidth,colspec={|Q[l,m,1.9cm]|X[0.85,l,m]|X[1.45,l,m]|X[0.75,l,m]|X[1.15,l,m]|},hlines,vlines,hspan=minimal,rowsep=1pt,colsep=2pt,row{1,3}={bg=headgray,font=\bfseries}}
MATA KULIAH & KODE & BOBOT (SKS)/jam & SEMESTER & TGL. PENYUSUNAN \\
Sistem Pendukung Keputusan & RTI254001 & 2 SKS/4 jam & 4 & 6 Januari 2026 \\
OTORISASI & Dosen Pengembang RPS & Koordinator RMK & \SetCell[c=2]{l} Ka PRODI &  \\
 & Tim Dosen Mata Kuliah &  & \SetCell[c=2]{l}{\vspace{8mm}\par Dr. Ely Setyo Astuti, S.T., M.T.} &  \\
\SetCell[r=11]{l} Capaian Pembelajaran (CP) & \SetCell[c=4]{l,bg=headgray,font=\bfseries} Capaian Pembelajaran Lulusan Program Studi (CPL-Prodi) &  &  &  \\
 & CPL07 & \SetCell[c=3]{l} Mampu mengembangkan sistem komputasi cerdas dan melakukan analisis data berbasis kecerdasan buatan untuk pengambilan keputusan. &  &  \\
 & CPL10 & \SetCell[c=3]{l} Mampu menganalisis persoalan kompleks dengan wawasan multidisiplin untuk menghasilkan solusi di bidang informatika dan ilmu komputer. &  &  \\
 & \SetCell[c=4]{l,bg=headgray,font=\bfseries} Capaian Pembelajaran mata kuliah (CPL-MK) &  &  &  \\
 & CPMK0706 & \SetCell[c=3]{l} Mahasiswa mampu mengimplementasikan berbagai metode MCDM dan sistem pendukung keputusan untuk menyelesaikan persoalan pengambilan keputusan multi-kriteria secara akurat. &  &  \\
 & CPMK1005 & \SetCell[c=3]{l} Mahasiswa mampu mengevaluasi dan membandingkan kinerja metode-metode SPK berdasarkan akurasi, konsistensi, dan interpretabilitas hasil untuk persoalan informatika. &  &  \\
 & \SetCell[c=4]{l,bg=headgray,font=\bfseries} Kemampuan akhir tiap tahapan belajar (Sub-CPMK) &  &  &  \\
 & SCPMK0706-02801 & \SetCell[c=3]{l} Mahasiswa mampu menjelaskan konsep dan evolusi sistem pendukung keputusan, teori pengambilan keputusan, serta menerapkan metode baseline MCDM yaitu WSM dan WPM pada kasus sederhana. &  &  \\
 & SCPMK0706-02802 & \SetCell[c=3]{l} Mahasiswa mampu menerapkan metode AHP untuk menghitung pembobotan kriteria berbasis pendapat decision maker baik untuk kriteria benefit maupun cost. &  &  \\
 & SCPMK0706-02803 & \SetCell[c=3]{l} Mahasiswa mampu menerapkan metode ROC, EDAS, MARCOS, dan MEREC untuk perankingan alternatif dan pembobotan objektif dalam pengambilan keputusan multi-kriteria. &  &  \\
 & SCPMK0706-02804 & \SetCell[c=3]{l} Mahasiswa mampu mengimplementasikan Group Decision Support Systems (GDSS) dan metode fuzzy untuk menangani ketidakpastian dan penilaian linguistik dalam SPK. &  &  \\
 & SCPMK1005-02805 & \SetCell[c=3]{l} Mahasiswa mampu merancang, mengimplementasikan, dan mengevaluasi sistem DSS lengkap menggunakan kombinasi metode MCDM dengan analisis metrik kinerja. &  &  \\
\SetCell[r=2]{l} Penilaian & \SetCell[c=4]{l} *Nilai CPMK didapat dari agregasi nilai SUB-CPMK yang dibebankan &  &  &  \\
 & \SetCell[c=4]{l}{\tiny\begin{tblr}{width=\linewidth,colspec={|Q[c,m,0.1\linewidth]|Q[c,m,0.16\linewidth]|X[l,m]|X[l,m]|X[l,m]|X[l,m]|X[l,m]|X[l,m]|Q[c,m,0.08\linewidth]|},hlines,vlines,hspan=minimal,rowsep=1pt,colsep=0.7pt,row{1,2}={bg=headgray,font=\bfseries}}
Kode CPMK & Kode SCPMK & \SetCell[c=6]{c} Bobot per Bentuk Penilaian & & & & & & Total Bobot \\
 & & Tugas & Kuis 1 & UTS & Kuis 2 & PBL & UAS & \\
\SetCell[r=4]{c} CPMK0706 & SCPMK0706-02801 & 5.6\%: SPK, MCDM, WSM/WPM & 5\%: konsep SPK & 5\%: kasus MCDM dasar & -- & -- & -- & 15.6\% \\
 & SCPMK0706-02802 & 8.4\%: AHP benefit dan cost & -- & 8\%: analisis AHP multi-kriteria & -- & -- & -- & 16.4\% \\
 & SCPMK0706-02803 & 11.2\%: ROC, EDAS, MARCOS, MEREC & -- & 7\%: perankingan alternatif & 5\%: metode objektif & -- & -- & 23.2\% \\
 & SCPMK0706-02804 & 5.6\%: GDSS dan fuzzy & -- & -- & -- & 5.7\%: proyek SPK & -- & 11.3\% \\
\SetCell[r=1]{c} CPMK1005 & SCPMK1005-02805 & 2.8\%: evaluasi dan desain DSS & -- & -- & -- & 5.7\%: proyek SPK & 25\%: presentasi DSS & 33.5\% \\
\SetCell[c=8]{r}\textbf{Total Per Penilaian} & & & & & & & & \textbf{100\%} \\
\end{tblr}} &  &  &  \\
Diskripsi Singkat Mata Kuliah & \SetCell[c=4]{l} Mata kuliah Sistem Pendukung Keputusan membangun kemampuan mahasiswa dalam merancang dan mengimplementasikan sistem pendukung keputusan berbasis metode Multi-Criteria Decision Making (MCDM) meliputi WSM, WPM, AHP, ROC, EDAS, MARCOS, MEREC, metode fuzzy, dan GDSS, serta mengevaluasi kinerja sistem DSS yang dibangun. &  &  &  \\
Bahan Kajian / Materi Pembelajaran & \SetCell[c=4]{l}\begin{enumerate}\item Konsep SPK, teori keputusan, dan MCDM dasar (WSM, WPM)\item Pembobotan subjektif: AHP (benefit dan cost)\item Metode perankingan lanjutan: ROC, EDAS, MARCOS, MEREC\item GDSS dan metode fuzzy untuk ketidakpastian\item Perancangan, implementasi, dan evaluasi sistem DSS\end{enumerate} &  &  &  \\
\SetCell[r=2]{l} Pustaka & \SetCell[c=4]{l} \textbf{Utama:}\begin{enumerate}\item Turban, E., Sharda, R., Delen, D. 2011. Decision Support and Business Intelligence Systems. 9th ed. Prentice Hall.\item Saaty, T.L. 1980. The Analytic Hierarchy Process. McGraw-Hill.\item Kusumadewi, S., Hartati, S., Harjoko, A., Wardoyo, R. 2006. Fuzzy Multi-Attribute Decision Making (Fuzzy MADM). Graha Ilmu.\end{enumerate} &  &  &  \\
 & \SetCell[c=4]{l} \textbf{Pendukung:} Materi LMS, modul praktikum Excel/Python, dan jurnal terkait metode MCDM dan DSS. &  &  &  \\
Media Pembelajaran & \SetCell[c=2]{l} \textbf{Software:} Microsoft Excel, Python/R (opsional), LMS. &  & \SetCell[c=2]{l} \textbf{Hardware:} Komputer/laptop, proyektor. &  \\
Nama Dosen Pengampu & \SetCell[c=4]{l} Tim Dosen Mata Kuliah &  &  &  \\
Matakuliah Syarat & \SetCell[c=4]{l} - &  &  &  \\
\end{longtblr}
\begin{longtblr}{width=\textwidth,colspec={|Q[c,m,0.06\textwidth]|X[1.35,l,m]|X[1.08,l,m]|X[1.16,l,m]|X[0.7,l,m]|X[1.24,l,m]|X[1.06,l,m]|X[0.98,l,m]|Q[c,m,0.05\textwidth]|},hlines,vlines,hspan=minimal,rowhead=2,row{1,2}={bg=headgreen,font=\bfseries},rowsep=1pt,colsep=1.5pt}
Mgg.\par Ke & Kemampuan Akhir Yang Direncanakan (Sub-CP-MK) & Bahan kajian (Materi Pembelajaran) & Bentuk dan Metode Pembelajaran & Estimasi Waktu & Pengalaman Belajar Mahasiswa & Kriteria \& Bentuk Penilaian & Indikator Penilaian & Bobot Penilaian (\%) \\
(1) & (2) & (3) & (4) & (5) & (6) & (7) & (8) & (9) \\
1 & SCPMK0706-02801 & Pengantar SPK, evolusi historis, teori pengambilan keputusan. & Modalitas: Blended Learning. Bentuk: Luring/Daring. Metode: Problem-based learning, studi kasus, dan presentasi. & 1 x 2 x 50' tatap muka; 1 x 2 x 50' tugas/praktik mandiri. & Mahasiswa mengerjakan aktivitas terarah untuk pengantar spk dan menunjukkan evidence hasil belajar. & Bentuk penilaian: Pengantar SPK. Mengacu rubrik RTM. & Ketepatan konsep; kualitas artefak/praktik; validasi dan dokumentasi. & 2.8 \\
2 & SCPMK0706-02801 & Konsep MCDM, elemen utama SPK, WSM dan WPM. & Modalitas: Blended Learning. Bentuk: Luring/Daring. Metode: Problem-based learning, studi kasus, dan presentasi. & 1 x 2 x 50' tatap muka; 1 x 2 x 50' tugas/praktik mandiri. & Mahasiswa mengerjakan aktivitas terarah untuk konsep mcdm dan wsm wpm dan menunjukkan evidence hasil belajar. & Bentuk penilaian: Konsep MCDM, WSM dan WPM. Mengacu rubrik RTM. & Ketepatan konsep; kualitas artefak/praktik; validasi dan dokumentasi. & 2.8 \\
3 & SCPMK0706-02802 & AHP 1: pembobotan subjektif kriteria benefit. & Modalitas: Blended Learning. Bentuk: Luring/Daring. Metode: Problem-based learning, studi kasus, dan presentasi. & 1 x 2 x 50' tatap muka; 1 x 2 x 50' tugas/praktik mandiri. & Mahasiswa mengerjakan aktivitas terarah untuk ahp pembobotan kriteria benefit dan menunjukkan evidence hasil belajar. & Bentuk penilaian: AHP 1 benefit. Mengacu rubrik RTM. & Ketepatan konsep; kualitas artefak/praktik; validasi dan dokumentasi. & 2.8 \\
4 & SCPMK0706-02802 & AHP 2: pembobotan subjektif kriteria cost-benefit. & Modalitas: Blended Learning. Bentuk: Luring/Daring. Metode: Problem-based learning, studi kasus, dan presentasi. & 1 x 2 x 50' tatap muka; 1 x 2 x 50' tugas/praktik mandiri. & Mahasiswa mengerjakan aktivitas terarah untuk ahp pembobotan kriteria cost-benefit dan menunjukkan evidence hasil belajar. & Bentuk penilaian: AHP 2 cost-benefit. Mengacu rubrik RTM. & Ketepatan konsep; kualitas artefak/praktik; validasi dan dokumentasi. & 2.8 \\
5 & SCPMK0706-02801 & Kuis 1: konsep SPK, MCDM dasar, WSM dan WPM. & Modalitas: Blended Learning. Bentuk: Luring/Daring. Metode: Problem-based learning, studi kasus, dan presentasi. & 1 x 2 x 50' tatap muka; 1 x 2 x 50' tugas/praktik mandiri. & Mahasiswa mengerjakan aktivitas terarah untuk kuis 1 konsep spk dan menunjukkan evidence hasil belajar. & Bentuk penilaian: Kuis 1 konsep SPK. Mengacu rubrik RTM. & Ketepatan konsep; kualitas artefak/praktik; validasi dan dokumentasi. & 5 \\
6 & SCPMK0706-02803 & ROC (Rank Order Centroid) dan EDAS. & Modalitas: Blended Learning. Bentuk: Luring/Daring. Metode: Problem-based learning, studi kasus, dan presentasi. & 1 x 2 x 50' tatap muka; 1 x 2 x 50' tugas/praktik mandiri. & Mahasiswa mengerjakan aktivitas terarah untuk roc dan edas dan menunjukkan evidence hasil belajar. & Bentuk penilaian: ROC dan EDAS. Mengacu rubrik RTM. & Ketepatan konsep; kualitas artefak/praktik; validasi dan dokumentasi. & 2.8 \\
7 & SCPMK0706-02803 & Metode MARCOS. & Modalitas: Blended Learning. Bentuk: Luring/Daring. Metode: Problem-based learning, studi kasus, dan presentasi. & 1 x 2 x 50' tatap muka; 1 x 2 x 50' tugas/praktik mandiri. & Mahasiswa mengerjakan aktivitas terarah untuk metode marcos dan menunjukkan evidence hasil belajar. & Bentuk penilaian: Metode MARCOS. Mengacu rubrik RTM. & Ketepatan konsep; kualitas artefak/praktik; validasi dan dokumentasi. & 2.8 \\
8 & SCPMK0706-02801, SCPMK0706-02802, SCPMK0706-02803 & UTS: analisis metode MCDM. & Modalitas: Blended Learning. Bentuk: Luring/Daring. Metode: Problem-based learning, studi kasus, dan presentasi. & 1 x 2 x 50' tatap muka; 1 x 2 x 50' tugas/praktik mandiri. & Mahasiswa mengerjakan aktivitas terarah untuk uts analisis metode mcdm dan menunjukkan evidence hasil belajar. & Bentuk penilaian: UTS analisis metode MCDM. Mengacu rubrik RTM. & Ketepatan konsep; kualitas artefak/praktik; validasi dan dokumentasi. & 20 \\
9 & SCPMK0706-02803 & MEREC: pembobotan objektif. & Modalitas: Blended Learning. Bentuk: Luring/Daring. Metode: Problem-based learning, studi kasus, dan presentasi. & 1 x 2 x 50' tatap muka; 1 x 2 x 50' tugas/praktik mandiri. & Mahasiswa mengerjakan aktivitas terarah untuk merec pembobotan objektif dan menunjukkan evidence hasil belajar. & Bentuk penilaian: MEREC pembobotan objektif. Mengacu rubrik RTM. & Ketepatan konsep; kualitas artefak/praktik; validasi dan dokumentasi. & 2.8 \\
10 & SCPMK0706-02803 & Group Decision Support Systems (GDSS). & Modalitas: Blended Learning. Bentuk: Luring/Daring. Metode: Problem-based learning, studi kasus, dan presentasi. & 1 x 2 x 50' tatap muka; 1 x 2 x 50' tugas/praktik mandiri. & Mahasiswa mengerjakan aktivitas terarah untuk group decision support systems dan menunjukkan evidence hasil belajar. & Bentuk penilaian: GDSS. Mengacu rubrik RTM. & Ketepatan konsep; kualitas artefak/praktik; validasi dan dokumentasi. & 2.8 \\
11 & SCPMK0706-02804 & Fuzzy DSS 1: ketidakpastian dan penilaian linguistik. & Modalitas: Blended Learning. Bentuk: Luring/Daring. Metode: Problem-based learning, studi kasus, dan presentasi. & 1 x 2 x 50' tatap muka; 1 x 2 x 50' tugas/praktik mandiri. & Mahasiswa mengerjakan aktivitas terarah untuk fuzzy dss ketidakpastian dan penilaian linguistik dan menunjukkan evidence hasil belajar. & Bentuk penilaian: Fuzzy DSS 1. Mengacu rubrik RTM. & Ketepatan konsep; kualitas artefak/praktik; validasi dan dokumentasi. & 2.8 \\
12 & SCPMK0706-02804 & Fuzzy DSS 2: implementasi fuzzy MCDM. & Modalitas: Blended Learning. Bentuk: Luring/Daring. Metode: Problem-based learning, studi kasus, dan presentasi. & 1 x 2 x 50' tatap muka; 1 x 2 x 50' tugas/praktik mandiri. & Mahasiswa mengerjakan aktivitas terarah untuk fuzzy dss implementasi fuzzy mcdm dan menunjukkan evidence hasil belajar. & Bentuk penilaian: Fuzzy DSS 2. Mengacu rubrik RTM. & Ketepatan konsep; kualitas artefak/praktik; validasi dan dokumentasi. & 2.8 \\
13 & SCPMK0706-02803 & Kuis 2: metode perankingan dan pembobotan objektif. & Modalitas: Blended Learning. Bentuk: Luring/Daring. Metode: Problem-based learning, studi kasus, dan presentasi. & 1 x 2 x 50' tatap muka; 1 x 2 x 50' tugas/praktik mandiri. & Mahasiswa mengerjakan aktivitas terarah untuk kuis 2 metode perankingan dan menunjukkan evidence hasil belajar. & Bentuk penilaian: Kuis 2 metode lanjutan. Mengacu rubrik RTM. & Ketepatan konsep; kualitas artefak/praktik; validasi dan dokumentasi. & 5 \\
14 & SCPMK1005-02805 & Metrik evaluasi SPK. & Modalitas: Blended Learning. Bentuk: Luring/Daring. Metode: Problem-based learning, studi kasus, dan presentasi. & 1 x 2 x 50' tatap muka; 1 x 2 x 50' tugas/praktik mandiri. & Mahasiswa mengerjakan aktivitas terarah untuk metrik evaluasi spk dan menunjukkan evidence hasil belajar. & Bentuk penilaian: Metrik evaluasi SPK. Mengacu rubrik RTM. & Ketepatan konsep; kualitas artefak/praktik; validasi dan dokumentasi. & 2.8 \\
15 & SCPMK1005-02805 & Perancangan sistem DSS dengan metode terkini. & Modalitas: Blended Learning. Bentuk: Luring/Daring. Metode: Problem-based learning, studi kasus, dan presentasi. & 1 x 2 x 50' tatap muka; 1 x 2 x 50' tugas/praktik mandiri. & Mahasiswa mengerjakan aktivitas terarah untuk perancangan sistem dss dan menunjukkan evidence hasil belajar. & Bentuk penilaian: Perancangan sistem DSS. Mengacu rubrik RTM. & Ketepatan konsep; kualitas artefak/praktik; validasi dan dokumentasi. & 2.8 \\
16 & SCPMK0706-02804, SCPMK1005-02805 & PBL: proyek implementasi DSS. & Modalitas: Blended Learning. Bentuk: Luring/Daring. Metode: Problem-based learning, studi kasus, dan presentasi. & 1 x 2 x 50' tatap muka; 1 x 2 x 50' tugas/praktik mandiri. & Mahasiswa mengerjakan aktivitas terarah untuk pbl proyek implementasi dss dan menunjukkan evidence hasil belajar. & Bentuk penilaian: PBL proyek implementasi DSS. Mengacu rubrik RTM. & Ketepatan konsep; kualitas artefak/praktik; validasi dan dokumentasi. & 11.4 \\
17 & SCPMK0706-02804, SCPMK1005-02805 & UAS: evaluasi dan presentasi proyek DSS. & Modalitas: Blended Learning. Bentuk: Luring/Daring. Metode: Problem-based learning, studi kasus, dan presentasi. & 1 x 2 x 50' tatap muka; 1 x 2 x 50' tugas/praktik mandiri. & Mahasiswa mengerjakan aktivitas terarah untuk uas evaluasi dan presentasi proyek dss dan menunjukkan evidence hasil belajar. & Bentuk penilaian: UAS evaluasi dan presentasi proyek DSS. Mengacu rubrik RTM. & Ketepatan konsep; kualitas artefak/praktik; validasi dan dokumentasi. & 25 \\
\end{longtblr}
